% Nguồn SGK Toán 9 Tập một — Bài 2, trang 11--18:
% https://cdn3.olm.vn/upload/thuvienso/4491668113-page-11-1773022924356.png
% https://cdn3.olm.vn/upload/thuvienso/4491668113-page-12-1773022924639.png
% https://cdn3.olm.vn/upload/thuvienso/4491668113-page-13-1773022925083.png
% https://cdn3.olm.vn/upload/thuvienso/4491668113-page-14-1773022925393.png
% https://cdn3.olm.vn/upload/thuvienso/4491668113-page-15-1773022925770.png
% https://cdn3.olm.vn/upload/thuvienso/4491668113-page-16-1773022926078.png
% https://cdn3.olm.vn/upload/thuvienso/4491668113-page-17-1773022926425.png
% https://cdn3.olm.vn/upload/thuvienso/4491668113-page-18-1773022926771.png
% Nguồn SBT: SBT TOAN 9 KNTT TAP 1.pdf, trang 9--13.
% SHA-256: d7ffd04618456682795977e5c47de7874161a35e8606e5bac1906881603a72f7

\section{Giải hệ hai phương trình bậc nhất hai ẩn}

\begin{lesson}
    \subsection{Kiến thức trọng tâm}

    \textbf{Khái niệm, thuật ngữ}
    \begin{itemize}
        \item Phương pháp thế
        \item Phương pháp cộng đại số
    \end{itemize}

    \textbf{Kiến thức, kĩ năng}

    Giải hệ hai phương trình bậc nhất hai ẩn bằng phương pháp thế, phương pháp cộng đại số hoặc sử dụng máy tính cầm tay.

    Một mảnh vườn được đánh thành nhiều luống, mỗi luống trồng cùng một số cây cải bắp. Hãy tính số cây cải bắp được trồng trên mảnh vườn đó, biết rằng:
    \begin{itemize}
        \item Nếu tăng thêm 8 luống, nhưng mỗi luống trồng ít đi 3 cây cải bắp thì số cải bắp của cả vườn sẽ ít đi 108 cây;
        \item Nếu giảm đi 4 luống, nhưng mỗi luống trồng tăng thêm 2 cây thì số cải bắp cả vườn sẽ tăng thêm 64 cây.
    \end{itemize}

    \subsubsection{Phương pháp thế}

    \textbf{Làm quen với phương pháp thế}

    \textbf{HĐ1} Cho hệ phương trình
    \[
        \begin{cases}
            x+y=3\\
            2x-3y=1.
        \end{cases}
    \]
    Giải hệ phương trình theo hướng dẫn sau:
    \begin{enumerate}
        \item Từ phương trình thứ nhất, biểu diễn $y$ theo $x$ rồi thế vào phương trình thứ hai để được một phương trình với một ẩn $x$. Giải phương trình một ẩn đó để tìm giá trị của $x$.
        \item Sử dụng giá trị tìm được của $x$ để tìm giá trị của $y$ rồi viết nghiệm của hệ phương trình đã cho.
    \end{enumerate}

    \answerline
    \answerline
    \answerline

    \begin{keyconcept}
        \textbf{Cách giải hệ phương trình bằng phương pháp thế:}
        \begin{itemize}
            \item \textbf{Bước 1.} Từ một phương trình của hệ, biểu diễn một ẩn theo ẩn kia rồi thế vào phương trình còn lại của hệ để được phương trình chỉ còn chứa một ẩn.
            \item \textbf{Bước 2.} Giải phương trình một ẩn vừa nhận được, từ đó suy ra nghiệm của hệ đã cho.
        \end{itemize}
    \end{keyconcept}

    \textbf{Ví dụ 1} Giải hệ phương trình
    \[
        \begin{cases}
            2x-y=3\\
            x+2y=4
        \end{cases}
    \]
    bằng phương pháp thế.

    \textbf{Giải}

    Từ phương trình thứ nhất của hệ ta có $y=2x-3$. Thế vào phương trình thứ hai của hệ, ta được $x+2(2x-3)=4$ hay $5x-6=4$, suy ra $x=2$.

    Từ đó $y=2\cdot2-3=1$. Vậy hệ phương trình đã cho có nghiệm là $(2;1)$.

    \textbf{Luyện tập 1} Giải các hệ phương trình sau bằng phương pháp thế:
    \begin{enumerate}[label=\alph*)]
        \item $\begin{cases}x-3y=2\\-2x+5y=1;\end{cases}$
        \item $\begin{cases}4x+y=-1\\7x+2y=-3.\end{cases}$
    \end{enumerate}

    \answerline
    \answerline
    \answerline

    Tuỳ theo hệ phương trình, ta có thể lựa chọn cách biểu diễn $x$ theo $y$ hoặc biểu diễn $y$ theo $x$.

    \textbf{Ví dụ 2} Giải hệ phương trình
    \[
        \begin{cases}
            x-y=-2\\
            2x-2y=8
        \end{cases}
    \]
    bằng phương pháp thế.

    \textbf{Giải}

    Từ phương trình thứ nhất ta có $x=y-2$. Thế vào phương trình thứ hai, ta được $2(y-2)-2y=8$ hay $0y-4=8$. \hfill (1)

    Do không có giá trị nào của $y$ thoả mãn hệ thức (1) nên hệ phương trình đã cho vô nghiệm.

    \textbf{Luyện tập 2} Giải hệ phương trình
    \[
        \begin{cases}
            -2x+y=3\\
            4x-2y=-4
        \end{cases}
    \]
    bằng phương pháp thế.

    \answerline
    \answerline
    \answerline

    \textbf{Ví dụ 3} Giải hệ phương trình
    \[
        \begin{cases}
            -x+y=-2\\
            3x-3y=6
        \end{cases}
    \]
    bằng phương pháp thế.

    \textbf{Giải}

    Từ phương trình thứ nhất ta có $y=x-2$. \hfill (2)

    Thế vào phương trình thứ hai, ta được $3x-3(x-2)=6$ hay $0x=0$. \hfill (3)

    Ta thấy mọi giá trị của $x$ đều thoả mãn hệ thức (3).

    Với giá trị tuỳ ý của $x$, giá trị tương ứng của $y$ được tính bởi (2).

    Vậy hệ phương trình đã cho có nghiệm là $(x;x-2)$ với $x\in\mathbb{R}$ tuỳ ý.

    \textbf{Luyện tập 3} Giải hệ phương trình
    \[
        \begin{cases}
            x+3y=-1\\
            3x+9y=-3
        \end{cases}
    \]
    bằng phương pháp thế.

    \answerline
    \answerline
    \answerline

    \textbf{Vận dụng 1} Xét bài toán trong tình huống mở đầu. Gọi $x$ là số luống trong vườn, $y$ là số cây cải bắp trồng ở mỗi luống ($x,y\in\mathbb{N}^{*}$).
    \begin{enumerate}[label=\alph*)]
        \item Lập hệ phương trình đối với hai ẩn $x$, $y$.
        \item Giải hệ phương trình nhận được ở câu a để tìm câu trả lời cho bài toán.
    \end{enumerate}

    \answerline
    \answerline
    \answerline

    \subsubsection{Phương pháp cộng đại số}

    \textbf{Làm quen với phương pháp cộng đại số}

    \textbf{HĐ2} Cho hệ phương trình
    \[
        \begin{cases}
            2x+2y=3\\
            x-2y=6.
        \end{cases}
    \]
    Ta thấy hệ số của $y$ trong hai phương trình là hai số đối nhau (tổng của chúng bằng 0). Từ đặc điểm đó, hãy giải hệ phương trình đã cho theo hướng dẫn sau:
    \begin{enumerate}
        \item Cộng từng vế của hai phương trình trong hệ để được phương trình một ẩn $x$. Giải phương trình này để tìm $x$.
        \item Sử dụng giá trị $x$ tìm được, thay vào một trong hai phương trình của hệ để tìm giá trị của $y$ rồi viết nghiệm của hệ phương trình đã cho.
    \end{enumerate}

    \answerline
    \answerline
    \answerline

    \begin{keyconcept}
        \textbf{Cách giải hệ phương trình bằng phương pháp cộng đại số:}

        Để giải một hệ hai phương trình bậc nhất hai ẩn \textit{có hệ số của cùng một ẩn nào đó trong hai phương trình bằng nhau hoặc đối nhau}, ta có thể làm như sau:
        \begin{itemize}
            \item \textbf{Bước 1.} Cộng hay trừ từng vế của hai phương trình trong hệ để được phương trình chỉ còn chứa một ẩn.
            \item \textbf{Bước 2.} Giải phương trình một ẩn vừa nhận được, từ đó suy ra nghiệm của hệ phương trình đã cho.
        \end{itemize}
    \end{keyconcept}

    \textbf{Ví dụ 4} Giải hệ phương trình
    \[
        \begin{cases}
            -2x+5y=12\\
            2x+3y=4
        \end{cases}
    \]
    bằng phương pháp cộng đại số.

    \textbf{Giải}

    Cộng từng vế hai phương trình ta được $8y=16$, suy ra $y=2$.

    Thế $y=2$ vào phương trình thứ hai ta được $2x+3\cdot2=4$, hay $2x=-2$, suy ra $x=-1$.

    Vậy hệ phương trình đã cho có nghiệm là $(-1;2)$.

    \textbf{Ví dụ 5} Giải hệ phương trình
    \[
        \begin{cases}
            5x-7y=9\\
            5x-3y=1
        \end{cases}
    \]
    bằng phương pháp cộng đại số.

    \textbf{Giải}

    Trừ từng vế của hai phương trình, ta được $(5x-5x)+(-7y+3y)=9-1$ hay $-4y=8$, suy ra $y=-2$.

    Thế $y=-2$ vào phương trình thứ nhất, ta được $5x-7\cdot(-2)=9$ hay $5x+14=9$, suy ra $x=-1$.

    Vậy hệ phương trình đã cho có nghiệm là $(-1;-2)$.

    \textbf{Luyện tập 4} Giải các hệ phương trình sau bằng phương pháp cộng đại số:
    \begin{enumerate}[label=\alph*)]
        \item $\begin{cases}-4x+3y=0\\4x-5y=-8;\end{cases}$
        \item $\begin{cases}4x+3y=0\\x+3y=9.\end{cases}$
    \end{enumerate}

    \answerline
    \answerline
    \answerline

    \textbf{Chú ý.} Trường hợp trong hệ phương trình đã cho không có hai hệ số của cùng một ẩn bằng nhau hay đối nhau, ta có thể đưa về trường hợp đã xét bằng cách nhân hai vế của mỗi phương trình với một số thích hợp (khác 0).

    \textbf{Ví dụ 6} Giải hệ phương trình
    \[
        \begin{cases}
            3x+2y=7\\
            2x-3y=-4
        \end{cases}
    \]
    bằng phương pháp cộng đại số.

    \textbf{Giải}

    Nhân hai vế của phương trình thứ nhất với 3 và nhân hai vế của phương trình thứ hai với 2, ta được:
    \[
        \begin{cases}
            9x+6y=21\\
            4x-6y=-8.
        \end{cases}
    \]

    Cộng từng vế hai phương trình của hệ mới, ta được $13x=13$ hay $x=1$.

    Thế $x=1$ vào phương trình thứ nhất của hệ đã cho, ta có $3\cdot1+2y=7$, suy ra $y=2$.

    Vậy hệ phương trình đã cho có nghiệm là $(1;2)$.

    \textbf{Luyện tập 5} Giải hệ phương trình
    \[
        \begin{cases}
            4x+3y=6\\
            -5x+2y=4
        \end{cases}
    \]
    bằng phương pháp cộng đại số.

    \answerline
    \answerline
    \answerline

    \textbf{Ví dụ 7} Giải hệ phương trình
    \[
        \begin{cases}
            3x-5y=2\\
            -6x+10y=-4
        \end{cases}
    \]
    bằng phương pháp cộng đại số.

    \textbf{Giải}

    Chia hai vế của phương trình thứ hai cho 2, ta được hệ
    \[
        \begin{cases}
            3x-5y=2\\
            -3x+5y=-2.
        \end{cases}
    \]

    Cộng từng vế hai phương trình của hệ mới ta có $0x+0y=0$. Hệ thức này luôn thoả mãn với các giá trị tuỳ ý của $x$ và $y$.

    Với giá trị tuỳ ý của $x$, giá trị của $y$ được tính nhờ hệ thức $3x-5y=2$, suy ra $y=\dfrac{3}{5}x-\dfrac{2}{5}$.

    Vậy hệ phương trình đã cho có nghiệm là $\left(x;\dfrac{3}{5}x-\dfrac{2}{5}\right)$ với $x\in\mathbb{R}$.

    \textbf{Luyện tập 6} Bằng phương pháp cộng đại số, giải hệ phương trình
    \[
        \begin{cases}
            -0{,}5x+0{,}5y=1\\
            -2x+2y=8.
        \end{cases}
    \]

    \answerline
    \answerline
    \answerline

    \subsubsection{Sử dụng máy tính cầm tay để tìm nghiệm của hệ hai phương trình bậc nhất hai ẩn}

    \textbf{Cách tìm nghiệm của hệ hai phương trình bậc nhất hai ẩn bằng máy tính cầm tay}

    Muốn tìm nghiệm của hệ hai phương trình bậc nhất hai ẩn bằng máy tính cầm tay (MTCT), chúng ta cần sử dụng loại máy có chức năng này (thường có phím MODE).

    Trước hết ta phải viết hệ phương trình cần tìm nghiệm dưới dạng:
    \[
        \begin{cases}
            a_1x+b_1y=c_1\\
            a_2x+b_2y=c_2.
        \end{cases}
    \]

    Chẳng hạn để tìm nghiệm của hệ
    \[
        \begin{cases}
            2x+3y-4=0\\
            5x+6y-7=0
        \end{cases}
    \]
    ta viết nó dưới dạng
    \[
        \begin{cases}
            2x+3y=4\\
            5x+6y=7.
        \end{cases}
    \]

    Khi đó, ta có $a_1=2$, $b_1=3$, $c_1=4$; $a_2=5$, $b_2=6$ và $c_2=7$. Lần lượt thực hiện các bước sau (với máy tính thích hợp):

    \textbf{Bước 1.} Vào chức năng giải hệ hai phương trình bậc nhất hai ẩn bằng cách bấm các phím \texttt{MODE} \texttt{5} \texttt{1} (xem màn hình sau bước 1, con trỏ ở vị trí $a_1$).

    % TODO(HINH-NGUON): bo sung asset/crop dung tu SGK trang 15, man hinh sau buoc 1; khong redraw xap xi bang TikZ.

    \textbf{Bước 2.} Nhập các hệ số $a_1=2$, $b_1=3$, $c_1=4$; $a_2=5$, $b_2=6$ và $c_2=7$ bằng cách bấm: \texttt{2 = 3 = 4 = 5 = 6 = 7 =} (xem màn hình sau bước 2).

    % TODO(HINH-NGUON): bo sung asset/crop dung tu SGK trang 15, man hinh sau buoc 2; khong redraw xap xi bang TikZ.

    \textbf{Bước 3.} Đọc kết quả: Sau khi kết thúc bước 2, bấm phím \texttt{=}, màn hình cho $x=-1$; bấm tiếp phím \texttt{=}, màn hình cho $y=2$ (xem màn hình sau bước 3). Ta hiểu nghiệm của hệ phương trình là $(-1;2)$.

    % TODO(HINH-NGUON): bo sung asset/crop dung tu SGK trang 15, man hinh sau buoc 3; khong redraw xap xi bang TikZ.

    \textbf{Chú ý}
    \begin{enumerate}
        \item Muốn xoá số vừa mới nhập thì bấm phím \texttt{AC}; muốn thay đổi số đã nhập ở một vị trí nào đó thì di chuyển con trỏ đến vị trí đó rồi nhập số mới.
        \item Bấm phím $\triangle$ hay $\triangledown$ để chuyển đổi hiển thị các giá trị của $x$ và $y$ trong kết quả.
        \item Nếu máy báo “Infinite Sol” thì hệ phương trình đã cho có vô số nghiệm. Nếu máy báo “No-Solution” thì hệ phương trình đã cho vô nghiệm.
    \end{enumerate}

    \textbf{Thực hành} Dùng MTCT thích hợp để tìm nghiệm của các hệ phương trình sau:
    \begin{enumerate}[label=\alph*)]
        \item $\begin{cases}2x+3y=-4\\-3x-7y=13;\end{cases}$
        \item $\begin{cases}2x+3y=1\\-x-1{,}5y=1;\end{cases}$
        \item $\begin{cases}8x-2y-6=0\\4x-y-3=0.\end{cases}$
    \end{enumerate}

    \answerline
    \answerline
    \answerline

    \textbf{Vận dụng 2} Thực hiện lần lượt các yêu cầu sau để tính số mililít dung dịch acid HCl nồng độ 20\% và số mililít dung dịch acid HCl nồng độ 5\% cần dùng để pha chế 2 lít dung dịch acid HCl nồng độ 10\%.
    \begin{enumerate}[label=\alph*)]
        \item Gọi $x$ là số mililít dung dịch acid HCl nồng độ 20\%, $y$ là số mililít dung dịch acid HCl nồng độ 5\% cần lấy. Hãy biểu thị theo $x$ và $y$:
        \begin{itemize}
            \item Thể tích của dung dịch acid HCl 10\% nhận được sau khi trộn lẫn hai dung dịch acid ban đầu.
            \item Tổng số gam acid HCl nguyên chất có trong hai dung dịch acid này.
        \end{itemize}
        \item Sử dụng kết quả ở câu a, hãy lập một hệ hai phương trình bậc nhất với hai ẩn là $x$, $y$. Giải hệ phương trình này để tính số mililít cần lấy của mỗi dung dịch acid HCl ở trên.
    \end{enumerate}

    \answerline
    \answerline
    \answerline

    \textbf{EM CÓ BIẾT?} \textit{(Đọc thêm)}

    \textbf{Giải hệ hai phương trình bậc nhất hai ẩn bằng phương pháp hình học}

    \textbf{1. Đường thẳng $ax+by=c$}

    Ta đã biết một phương trình bậc nhất hai ẩn luôn có vô số nghiệm; mỗi nghiệm là một cặp số $(x_0;y_0)$; mỗi cặp số lại có thể xem là tọa độ của một điểm trên mặt phẳng tọa độ $Oxy$. Vậy trên mặt phẳng tọa độ, các điểm mà tọa độ của chúng là các nghiệm của một phương trình bậc nhất hai ẩn có quan hệ với nhau như thế nào?

    \textbf{Người ta đã chứng minh được rằng:} Trong mặt phẳng tọa độ $Oxy$, tập hợp tất cả các điểm có tọa độ thoả mãn phương trình bậc nhất hai ẩn $ax+by=c$ tạo nên một đường thẳng. Đường thẳng đó gọi là đường thẳng $ax+by=c$. Nếu kí hiệu đường thẳng đó là $\Delta$ thì ta viết $\Delta: ax+by=c$.
    \begin{itemize}
        \item Khi $a\ne0$ và $b\ne0$, đường thẳng $\Delta$ trùng với đồ thị hàm số $y=-\dfrac{a}{b}x+\dfrac{c}{b}$;
        \item Khi $a=0$ và $b\ne0$, phương trình $ax+by=c$ có thể đưa về dạng $y=m$ (với $m=\dfrac{c}{b}$); $\Delta$ là đường thẳng song song hoặc trùng với trục hoành và cắt trục tung tại điểm $(0;m)$;
        \item Khi $a\ne0$ và $b=0$, phương trình $ax+by=c$ có thể đưa về dạng $x=n$ (với $n=\dfrac{c}{a}$); $\Delta$ là đường thẳng song song hoặc trùng với trục tung và cắt trục hoành tại điểm $(n;0)$.
    \end{itemize}

    Dưới đây là một số ví dụ (Hình 1.3a, b, c):

    % TODO(HINH-NGUON): bo sung asset/crop dung tu SGK trang 17, Hinh 1.3a,b,c; khong redraw xap xi bang TikZ.

    \textbf{2. Giải hệ hai phương trình bậc nhất hai ẩn bằng hình học}

    Ta đã biết, mỗi nghiệm của hệ phương trình bậc nhất hai ẩn
    \[
        \begin{cases}
            ax+by=c\\
            a'x+b'y=c'
        \end{cases}
        \tag{*}
    \]
    là một nghiệm chung của hai phương trình trong hệ (*). Nghiệm chung ấy tương ứng với điểm chung của hai đường thẳng $\Delta: ax+by=c$ và $\Delta': a'x+b'y=c'$, tức là giao điểm của $\Delta$ và $\Delta'$. Do đó ta có thể giải hệ (*) bằng cách vẽ hai đường thẳng $\Delta$ và $\Delta'$ rồi tìm tọa độ điểm chung của chúng.

    Từ đó, ta thấy chỉ có thể xảy ra 3 trường hợp:
    \begin{enumerate}
        \item $\Delta$ và $\Delta'$ cắt nhau (có một điểm chung): Hệ (*) có một nghiệm duy nhất.
        \item $\Delta$ và $\Delta'$ song song với nhau (không có điểm chung): Hệ (*) vô nghiệm.
        \item $\Delta$ và $\Delta'$ trùng nhau (mỗi điểm của $\Delta$ đều là điểm chung): Hệ (*) có vô số nghiệm.
    \end{enumerate}

    \textbf{Ví dụ 1} Để giải hệ phương trình
    \[
        \begin{cases}
            2x-y=3\\
            x+2y=4
        \end{cases}
    \]
    ta làm như sau (H.1.4):
    \begin{itemize}
        \item Vẽ đường thẳng $2x-y=3$ qua hai điểm $A(0;-3)$ và $B(1;-1)$.
        \item Vẽ đường thẳng $x+2y=4$ qua hai điểm $C(0;2)$ và $D(4;0)$.
    \end{itemize}

    Hai đường thẳng cắt nhau tại $M(2;1)$ nên $(2;1)$ là nghiệm duy nhất của hệ đã cho.

    % TODO(HINH-NGUON): bo sung asset/crop dung tu SGK trang 18, Hinh 1.4; khong redraw xap xi bang TikZ.

    \textbf{Ví dụ 2} Xét hệ phương trình
    \[
        \begin{cases}
            -x+y=2\\
            -2x+2y=8
        \end{cases}
    \]
    ta thấy đường thẳng $\Delta: -x+y=2$ là đồ thị của hàm số $y=x+2$; đường thẳng $\Delta': -2x+2y=8$ là đồ thị của hàm số $y=x+4$ (H.1.5).

    Hai đường thẳng $\Delta$ và $\Delta'$ phân biệt và có cùng hệ số góc là 1 nên chúng song song với nhau.

    Do đó, hệ phương trình đã cho vô nghiệm.

    % TODO(HINH-NGUON): bo sung asset/crop dung tu SGK trang 18, Hinh 1.5; khong redraw xap xi bang TikZ.

    \textbf{Ví dụ 3} Xét hệ phương trình
    \[
        \begin{cases}
            -x+y=2\\
            -2x+2y=4
        \end{cases}
    \]
    ta thấy hai đường thẳng $-x+y=2$ và $-2x+2y=4$ cùng là đồ thị của hàm số $y=x+2$ nên chúng trùng nhau. Vậy hệ phương trình đã cho có vô số nghiệm.
\end{lesson}

\begin{exercises}
    \subsection{Bài tập}

    \begin{questions}[resume=SGK]
        \item Giải các hệ phương trình sau bằng phương pháp thế:
        \begin{questions}
            \item $\begin{cases}x-y=3\\3x-4y=2\end{cases};$

            \answerline
            \answerline
            \answerline

            \item $\begin{cases}7x-3y=13\\4x+y=2\end{cases};$

            \answerline
            \answerline
            \answerline

            \item $\begin{cases}0{,}5x-1{,}5y=1\\-x+3y=2\end{cases}.$

            \answerline
            \answerline
            \answerline
        \end{questions}

        \item Giải các hệ phương trình sau bằng phương pháp cộng đại số:
        \begin{questions}
            \item $\begin{cases}3x+2y=6\\2x-2y=14\end{cases};$

            \answerline
            \answerline
            \answerline

            \item $\begin{cases}0{,}3x+0{,}5y=3\\1{,}5x-2y=1{,}5\end{cases};$

            \answerline
            \answerline
            \answerline

            \item $\begin{cases}-2x+6y=8\\3x-9y=-12\end{cases}.$

            \answerline
            \answerline
            \answerline
        \end{questions}

        \item Cho hệ phương trình
        \[
            \begin{cases}
                2x-y=-3\\
                -2m^2x+9y=3(m+3),
            \end{cases}
        \]
        trong đó $m$ là số đã cho. Giải hệ phương trình trong mỗi trường hợp sau:
        \begin{questions}
            \item $m=-2$;

            \answerline
            \answerline

            \item $m=-3$;

            \answerline
            \answerline

            \item $m=3$.

            \answerline
            \answerline
        \end{questions}

        \item Dùng MTCT thích hợp để tìm nghiệm của các hệ phương trình sau:
        \begin{questions}
            \item $\begin{cases}12x-5y+24=0\\-5x-3y-10=0\end{cases};$

            \answerline

            \item $\begin{cases}\dfrac{1}{3}x-y=\dfrac{2}{3}\\x-3y=2\end{cases};$

            \answerline

            \item $\begin{cases}3x-2y=1\\-x+\dfrac{2}{3}y=0\end{cases};$

            \answerline

            \item $\begin{cases}\dfrac{4}{9}x-\dfrac{3}{5}y=11\\\dfrac{2}{9}x+\dfrac{1}{5}y=-2\end{cases}.$

            \answerline
        \end{questions}
    \end{questions}
\end{exercises}

\begin{practice}
    \subsection{Luyện tập}

    \begin{questions}[resume=SBT]
        \item Giải các hệ phương trình sau bằng phương pháp thế:
        \begin{questions}
            \item $\begin{cases}x+2y=8\\\dfrac{1}{2}x-y=18\end{cases};$

            \answerline
            \answerline
            \answerline

            \item $\begin{cases}0{,}2x+0{,}5y=0{,}7\\4x+10y=9\end{cases};$

            \answerline
            \answerline
            \answerline

            \item $\begin{cases}-2x+3y=1\\\dfrac{1}{3}x-\dfrac{1}{2}y=-\dfrac{1}{6}\end{cases}.$

            \answerline
            \answerline
            \answerline
        \end{questions}

        \item Giải các hệ phương trình sau bằng phương pháp cộng đại số:
        \begin{questions}
            \item $\begin{cases}3x-7y=-14\\5x+2y=45\end{cases};$

            \answerline
            \answerline
            \answerline

            \item $\begin{cases}x-0{,}5y=-3\\2x-y=6\end{cases};$

            \answerline
            \answerline
            \answerline

            \item $\begin{cases}2x+3y=3\\\dfrac{2}{3}x+y=1\end{cases}.$

            \answerline
            \answerline
            \answerline
        \end{questions}

        \item Dùng MTCT, tìm nghiệm của các hệ phương trình sau:
        \begin{questions}
            \item $\begin{cases}\sqrt{3}x+3y=1\\2x-\sqrt{3}y=\sqrt{3}\end{cases};$

            \answerline

            \item $\begin{cases}2{,}5x-3{,}5y=0{,}5\\-0{,}5x+0{,}7y=1\end{cases};$

            \answerline

            \item $\begin{cases}\dfrac{x}{5}+\dfrac{y}{2}=5\\0{,}4x+y=1\end{cases}.$

            \answerline
        \end{questions}

        \item Tìm $a$ và $b$ để đường thẳng $y=ax+b$ đi qua hai điểm $(4;1)$ và $(-4;-3)$.

        \answerline
        \answerline

        \item Cho hệ phương trình
        \[
            \begin{cases}
                x+3y=1\\
                2x+my=5.
            \end{cases}
        \]
        \begin{questions}
            \item Giải hệ với $m=1$.

            \answerline
            \answerline

            \item Chứng tỏ rằng hệ đã cho vô nghiệm khi $m=6$.

            \answerline
            \answerline
        \end{questions}

        \item Trong kinh tế học, đường IS là một phương trình bậc nhất biểu diễn tất cả các kết hợp thu nhập $Y$ và lãi suất $r$ để duy trì trạng thái cân bằng của thị trường hàng hóa trong nền kinh tế. Đường LM là một phương trình bậc nhất biểu diễn tất cả các kết hợp giữa thu nhập $Y$ và lãi suất $r$ duy trì trạng thái cân bằng của thị trường tiền tệ. Trong một nền kinh tế, giả sử mức thu nhập cân bằng $Y$ (tính bằng triệu đô la) và lãi suất cân bằng $r$ thỏa mãn hệ phương trình
        \[
            \begin{cases}
                0{,}06Y-5\,000r=240\\
                0{,}06Y+6\,000r=900.
            \end{cases}
        \]
        Tìm mức thu nhập và lãi suất cân bằng.

        \answerline
        \answerline
        \answerline

        \item Phương trình cung và phương trình cầu của một loại thiết bị kĩ thuật số cá nhân mới là:

        Phương trình cầu: $p=150-0{,}00001x$;

        Phương trình cung: $p=60+0{,}00002x$;

        trong đó $p$ là giá mỗi đơn vị (tính bằng đô la) và $x$ là số lượng đơn vị sản phẩm. Tìm \textit{điểm cân bằng} của thị trường này, tức là điểm $(p;x)$ thỏa mãn cả hai phương trình cung và cầu.

        \answerline
        \answerline
        \answerline

        \item Thầy Nam dạy Toán đang thiết kế một bài kiểm tra trắc nghiệm gồm hai loại câu hỏi, câu hỏi đúng/sai và câu hỏi nhiều lựa chọn. Bài kiểm tra sẽ được tính thang điểm 100, trong đó mỗi câu hỏi đúng/sai có giá trị 2 điểm và mỗi câu hỏi nhiều lựa chọn có giá trị 4 điểm. Thầy Nam muốn số câu hỏi nhiều lựa chọn gấp đôi số câu hỏi đúng/sai.
        \begin{questions}
            \item Gọi số câu hỏi đúng/sai là $x$, số câu hỏi nhiều lựa chọn là $y$ ($x,y\in\mathbb{N}^{*}$). Viết hệ hai phương trình biểu thị số lượng của từng loại câu hỏi.

            \answerline
            \answerline

            \item Giải hệ phương trình trong câu a để biết số lượng câu hỏi mỗi loại trong bài kiểm tra là bao nhiêu.

            \answerline
            \answerline
        \end{questions}
    \end{questions}
\end{practice}